\documentclass[12pt, a4paper]{article}

\usepackage[utf8]{inputenc}
\usepackage[T1]{fontenc}
\usepackage[french]{babel}
\usepackage{amsmath, amssymb, amsthm}
\usepackage{bm}
\usepackage{graphicx}
\usepackage{hyperref}
\usepackage{geometry}
\usepackage{booktabs}
\usepackage{xcolor}
\usepackage{mathtools}

\geometry{margin=2.5cm}

\hypersetup{
    colorlinks=true,
    linkcolor=blue!70!black,
    citecolor=blue!70!black,
}

\newtheorem{proposition}{Proposition}
\newtheorem{remarque}{Remarque}

\title{%
    \textbf{Représentation des HRTFs par Harmoniques Sphériques}\\[0.4em]
    \large Formalisme mathématique et justification pour l'encodage cross-modal
}
\author{Mehdi Maqlach}
\date{\today}

% -----------------------------------------------------------------------
\begin{document}
\maketitle
\tableofcontents
\vspace{1em}

% -----------------------------------------------------------------------
\section{Introduction}

Une \textit{Head-Related Transfer Function} (HRTF) caractérise la façon dont le son
est filtré par la géométrie du corps humain — pavillon de l'oreille, tête, torse —
avant d'atteindre le canal auditif. Formellement, pour un sujet donné, la HRTF est
une fonction

\begin{equation}
    H : S^2 \times \mathbb{R}^+ \;\longrightarrow\; \mathbb{C}
    \qquad
    (\theta, \varphi, f) \;\longmapsto\; H(\theta, \varphi, f)
\end{equation}

où $(\theta, \varphi) \in S^2$ désigne la direction de la source sonore (azimut et
élévation sur la sphère unité) et $f$ la fréquence.

\medskip

Dans les bases de données HRTF mesurées (e.g.\ HUTUBS, SONICOM), cette fonction
est échantillonnée sur un ensemble fini de directions
$\{(\theta_i, \varphi_i)\}_{i=1}^{N}$, dont le cardinal $N$ varie selon le protocole
de mesure~: $N = 2114$, $2818$ ou $8802$ dans la base \textit{Database-Master V2-1}.
Cette hétérogénéité empêche de représenter directement les HRTFs comme un tableau
numpy de forme fixe, ce qui est un prérequis pour tout pipeline d'apprentissage
automatique.

\medskip

Ce document présente la \textbf{décomposition en harmoniques sphériques} (SH) comme
solution rigoureuse à ce problème : projeter la HRTF sur une base orthonormée de
fonctions définies sur $S^2$, indépendante de la grille de mesure. La représentation
obtenue est de taille fixe, compacte, et physiquement motivée.

% -----------------------------------------------------------------------
\section{Harmoniques sphériques}

\subsection{Définition}

Les harmoniques sphériques réelles d'ordre $l \geq 0$ et de degré $|m| \leq l$
sont définies par~:

\begin{equation}
    Y_l^m(\theta, \varphi) =
    \begin{cases}
        \sqrt{2}\, K_l^m \cos(m\varphi)\, P_l^m(\cos\theta) & m > 0 \\
        K_l^0\, P_l^0(\cos\theta)                            & m = 0 \\
        \sqrt{2}\, K_l^{|m|} \sin(|m|\varphi)\, P_l^{|m|}(\cos\theta) & m < 0
    \end{cases}
\end{equation}

où $P_l^m$ sont les polynômes de Legendre associés et

\begin{equation}
    K_l^m = \sqrt{\frac{(2l+1)}{4\pi} \frac{(l-|m|)!}{(l+|m|)!}}
\end{equation}

est la constante de normalisation. Pour un ordre de troncature $L$, le nombre total
de fonctions de base est

\begin{equation}
    \boxed{N_{\mathrm{SH}}(L) = (L+1)^2}
\end{equation}

\subsection{Propriété fondamentale : orthonormalité}

Les harmoniques sphériques forment une base orthonormée de $L^2(S^2)$~:

\begin{equation}
    \int_{S^2} Y_l^m(\omega)\, Y_{l'}^{m'}(\omega)\, \mathrm{d}\omega
    = \delta_{ll'}\, \delta_{mm'}
    \label{eq:ortho}
\end{equation}

Cette propriété est l'équivalent sphérique de l'orthogonalité des harmoniques
de Fourier sur le cercle. Elle garantit que les coefficients de la décomposition
sont \emph{décorrélés} : chaque coefficient capture une composante spatiale
indépendante.

\subsection{Complétude}

Toute fonction de carré intégrable sur $S^2$ peut être approchée arbitrairement
bien par une combinaison finie d'harmoniques sphériques~:

\begin{equation}
    H(\theta, \varphi, f) \;=\; \lim_{L \to \infty} \sum_{l=0}^{L} \sum_{m=-l}^{l}
    a_l^m(f)\, Y_l^m(\theta, \varphi)
    \label{eq:decomp}
\end{equation}

En pratique, les HRTFs sont des fonctions régulières sur $S^2$ (pas de
discontinuité abrupte), ce qui assure une convergence rapide : un ordre $L$
modéré suffit à capturer l'essentiel de l'information spatiale.

% -----------------------------------------------------------------------
\section{Décomposition d'une HRTF}

\subsection{Problème discret}

En pratique, on ne dispose que de $N$ mesures aux positions
$\{\omega_i\}_{i=1}^{N}$.
Pour une fréquence $f$ et une oreille données, on cherche les coefficients
$\mathbf{a}(f) = [a_0^0(f),\, a_1^{-1}(f),\, a_1^0(f),\, \ldots,\,
a_L^L(f)]^\top \in \mathbb{R}^{(L+1)^2}$
qui minimisent l'erreur de reconstruction aux positions mesurées.

On définit la \textbf{matrice de design} SH :

\begin{equation}
    \mathbf{Y} \;\in\; \mathbb{R}^{N \times (L+1)^2},
    \qquad
    \mathbf{Y}_{i,\,lm} = Y_l^m(\omega_i)
\end{equation}

et le vecteur de mesures (log-magnitude en dB) :

\begin{equation}
    \mathbf{h}(f) = \bigl[\hat{H}(\omega_1, f),\; \hat{H}(\omega_2, f),\;
    \ldots,\; \hat{H}(\omega_N, f)\bigr]^\top \in \mathbb{R}^N
\end{equation}

Le système linéaire à résoudre est :

\begin{equation}
    \mathbf{Y}\, \mathbf{a}(f) = \mathbf{h}(f)
    \label{eq:system}
\end{equation}

\subsection{Solution par moindres carrés}

Lorsque $N > (L+1)^2$ (système sur-déterminé, ce qui est notre cas), la solution
aux moindres carrés est donnée par la pseudo-inverse de Moore-Penrose~:

\begin{equation}
    \boxed{
    \hat{\mathbf{a}}(f) = \mathbf{Y}^+ \mathbf{h}(f)
    = \bigl(\mathbf{Y}^\top \mathbf{Y}\bigr)^{-1} \mathbf{Y}^\top\, \mathbf{h}(f)
    }
    \label{eq:lstsq}
\end{equation}

\begin{remarque}
La matrice $\mathbf{Y}$ ne dépend \emph{que} des positions de mesure, pas des
données HRTF. Elle est donc calculée \textbf{une seule fois} par sujet et
réutilisée pour toutes les fréquences et les deux oreilles.
\end{remarque}

\subsection{Représentation finale}

En appliquant \eqref{eq:lstsq} à chaque fréquence $f_k$ ($k = 1, \ldots, N_f$)
et aux deux oreilles ($e \in \{L, R\}$), on obtient le tenseur de coefficients :

\begin{equation}
    \mathbf{A} \;\in\; \mathbb{R}^{(L+1)^2 \times N_f \times 2}
\end{equation}

Ce tenseur constitue la \textbf{représentation SH de la HRTF} du sujet.
Sa taille est \emph{entièrement déterminée par le choix de $L$}, indépendamment
du nombre $N$ de positions de mesure.

\medskip

\noindent Récapitulatif numérique pour notre base de données (FIR 256 taps,
$f_s = 44100$\,Hz, donc $N_f = 129$)~:

\begin{center}
\begin{tabular}{ccccc}
\toprule
Ordre $L$ & $(L+1)^2$ & Shape $\mathbf{A}$ & Taille (float32) & Ratio vs.\ $N=2114$ \\
\midrule
3  & 16  & $[16,\;129,\;2]$  & 16\,kB  & $\times 0.008$ \\
5  & 36  & $[36,\;129,\;2]$  & 37\,kB  & $\times 0.017$ \\
10 & 121 & $[121,\;129,\;2]$ & 124\,kB & $\times 0.058$ \\
20 & 441 & $[441,\;129,\;2]$ & 453\,kB & $\times 0.21$  \\
\bottomrule
\end{tabular}
\end{center}

% -----------------------------------------------------------------------
\section{Justification pour l'apprentissage cross-modal}

\subsection{Problème d'hétérogénéité résolu proprement}

La décomposition SH est \textbf{la seule approche} qui résout le problème de
grilles hétérogènes sans perte d'information et sans artefact~:

\begin{itemize}
    \item \textbf{Troncature} au minimum $N_\mathrm{min}$ : supprime des mesures
    (perte d'information irréversible, ordre de troncature arbitraire).
    \item \textbf{Moyenne spatiale} : réduction à une dimension, toute information
    directionnelle détruite.
    \item \textbf{Interpolation} sur grille commune : introduit des artefacts si
    la densité de la grille source est insuffisante.
    \item \textbf{Décomposition SH} : utilise \emph{toutes} les mesures disponibles
    (D1 profite de ses 8802 points), produit une représentation de taille fixe,
    et maximise la précision sous contrainte de compacité.
\end{itemize}

\subsection{Décorrélation des caractéristiques}

L'orthogonalité \eqref{eq:ortho} garantit que les coefficients
$a_l^m$ sont décorrélés entre eux : chaque coefficient capture une
\emph{composante indépendante} de la variation spatiale. Cette propriété est
favorable pour un encodeur de type MLP ou CNN, dont l'optimisation est facilitée
par des entrées non corrélées.

\subsection{Hiérarchie spatiale interprétable}

Les harmoniques sphériques de bas ordre capturent les variations \textbf{globales}
de la HRTF (différence gauche/droite, bas/haut), tandis que les ordres élevés
encodent les détails \textbf{fins} du pavillon. Cette hiérarchie naturelle est
cohérente avec l'objectif du modèle~: retrouver la signature spectrale d'une oreille
depuis sa photo, qui contient justement ces niveaux de détail emboîtés.

\subsection{Reconstruction parfaite et robustesse}

Depuis les coefficients $\mathbf{A}$, on peut reconstruire la HRTF à
\emph{n'importe quelle direction} $(\theta, \varphi)$ par~:

\begin{equation}
    \hat{H}(\theta, \varphi, f) = \sum_{l=0}^{L} \sum_{m=-l}^{l}
    a_l^m(f)\, Y_l^m(\theta, \varphi)
\end{equation}

Cette propriété est utile en inférence~: une fois le modèle entraîné, il prédit
$\mathbf{A}$ depuis la photo d'oreille, puis on synthétise la HRTF dans
n'importe quelle direction pour la spatialisation binaural.

\subsection{Lien avec l'acoustique physique}

Les harmoniques sphériques sont les \textbf{solutions propres de l'équation de
Laplace sur $S^2$}, et par extension les modes naturels de la propagation
acoustique en champ libre. La décomposition SH d'une HRTF est donc
\emph{physiquement motivée}~: elle sépare les contributions des différentes
échelles spatiales de diffraction sur la tête et le pavillon. C'est aussi la base
mathématique de l'Ambisonics, standard industriel de l'audio spatial.

% -----------------------------------------------------------------------
\section{Choix de l'ordre $L$}

\subsection{Compromis précision / compacité}

L'erreur de reconstruction décroît quand $L$ augmente. Pour les HRTFs, la
littérature~\cite{Evans98, Zotkin09} montre qu'un ordre $L = 5$ à $L = 10$
capture plus de 95\,\% de l'énergie spatiale.

\subsection{Condition de Nyquist sphérique}

Pour que la décomposition soit non-aliasée, il faut que le nombre de mesures
$N$ vérifie~:

\begin{equation}
    N \;\geq\; (L+1)^2
\end{equation}

Avec $N_\mathrm{min} = 2114$ positions, on peut monter jusqu'à $L \approx 44$
sans sous-détermination.

\subsection{Recommandation pratique}

Pour une première version du modèle cross-modal~:

\begin{center}
\begin{tabular}{ll}
\toprule
Contexte & Ordre recommandé \\
\midrule
Prototype rapide      & $L = 5$  $(36 \text{ coeff.})$ \\
Version de référence  & $L = 10$ $(121 \text{ coeff.})$ \\
Haute fidélité        & $L = 20$ $(441 \text{ coeff.})$ \\
\bottomrule
\end{tabular}
\end{center}

% -----------------------------------------------------------------------
\section{Implémentation}

En Python, la décomposition se résume à trois étapes~:

\begin{enumerate}
    \item Lire les positions $(\theta_i, \varphi_i)$ depuis le fichier SOFA.
    \item Construire $\mathbf{Y}$ avec \texttt{scipy.special.sph\_harm}.
    \item Résoudre $\mathbf{A} = \mathbf{Y}^+ \mathbf{H}$ avec
          \texttt{numpy.linalg.lstsq}.
\end{enumerate}

La matrice $\mathbf{Y}^+$ est calculée une seule fois par sujet (coût
$\mathcal{O}(N \cdot (L+1)^2)$), puis appliquée à toutes les fréquences
simultanément par produit matriciel~:

\begin{equation}
    \mathbf{A} = \mathbf{Y}^+ \mathbf{H}
    \qquad
    \mathbf{H} \in \mathbb{R}^{N \times N_f \times 2}
    \;\longrightarrow\;
    \mathbf{A} \in \mathbb{R}^{(L+1)^2 \times N_f \times 2}
\end{equation}

% -----------------------------------------------------------------------
\section{Conclusion}

La décomposition en harmoniques sphériques offre une représentation des HRTFs~:
\begin{itemize}
    \item \textbf{de taille fixe} — $(L+1)^2 \times N_f \times 2$, indépendante de
    la grille de mesure~;
    \item \textbf{sans perte} — toutes les mesures contribuent à la solution des
    moindres carrés~;
    \item \textbf{physiquement motivée} — base propre de la propagation acoustique
    sphérique~;
    \item \textbf{interprétable} — hiérarchie ordre bas / ordre élevé cohérente avec
    la géométrie du pavillon~;
    \item \textbf{réversible} — reconstruction à n'importe quelle direction par
    évaluation de la série tronquée.
\end{itemize}

Elle constitue donc le prétraitement le plus rigoureux pour alimenter un
HRTFEncoder dans une architecture de retrieval cross-modal oreille $\leftrightarrow$
HRTF.

% -----------------------------------------------------------------------
\begin{thebibliography}{9}

\bibitem{Evans98}
M.~A. Evans, J.~A. S. Angus, and A.~I. Tew,
\textit{Analyzing head-related transfer function measurements using surface spherical harmonics},
J.\ Acoust.\ Soc.\ Am., 104(4):2400--2411, 1998.

\bibitem{Zotkin09}
D.~N. Zotkin, R.~Duraiswami, and N.~A. Gumerov,
\textit{Regularized HRTF fitting using spherical harmonics},
in Proc.\ IEEE WASPAA, 2009.

\bibitem{Rafaely15}
B.~Rafaely,
\textit{Fundamentals of Spherical Array Processing},
Springer, 2015.

\bibitem{Chen19}
T.~Chen, S.\ Kornblith, M.\ Norouzi, and G.\ Hinton,
\textit{A simple framework for contrastive learning of visual representations (SimCLR)},
ICML, 2020.

\end{thebibliography}

\end{document}
